% --- CORRECCION: Forzar interpretacion raw de la entrada ---
\UseRawInputEncoding
\documentclass[aspectratio=169, 10pt, xcolor=table]{beamer}

% --- Configuracion del Tema ---
\usetheme{Madrid}
\usecolortheme{dolphin}

% --- Paquetes necesarios ---
\usepackage[utf8]{inputenc}
\usepackage[T1]{fontenc}
\usepackage[spanish,es-noquoting]{babel}
\usepackage{amsmath, amssymb}
\usepackage{mathtools}
\usepackage{booktabs}
\usepackage{graphicx}
\usepackage{listings}
\usepackage{tikz}
\usetikzlibrary{arrows.meta, positioning, shapes.geometric, babel}
\usepackage{etoolbox}
\usepackage{upquote}

% --- Estilo para codigo SQL, bash, YAML ---
\definecolor{codegreen}{rgb}{0,0.6,0}
\definecolor{codegray}{rgb}{0.5,0.5,0.5}
\definecolor{codepurple}{rgb}{0.58,0,0.82}
\definecolor{backcolour}{rgb}{0.95,0.95,0.92}
\lstdefinestyle{bashstyle}{
    language=bash,
    backgroundcolor=\color{backcolour},
    commentstyle=\color{codegreen},
    keywordstyle=\color{blue},
    numberstyle=\tiny\color{codegray},
    stringstyle=\color{codepurple},
    basicstyle=\ttfamily\scriptsize,
    breaklines=true,
    frame=single,
    showstringspaces=false,
    literate={á}{{\'a}}1 {é}{{\'e}}1 {í}{{\'i}}1 {ó}{{\'o}}1 {ú}{{\'u}}1 {ñ}{{\~n}}1 {ü}{{\"u}}1
}
\lstdefinestyle{sqlstyle}{
    language=SQL,
    backgroundcolor=\color{backcolour},
    commentstyle=\color{codegreen},
    keywordstyle=\color{blue},
    numberstyle=\tiny\color{codegray},
    stringstyle=\color{codepurple},
    basicstyle=\ttfamily\scriptsize,
    breaklines=true,
    frame=single,
    showstringspaces=false,
    literate={á}{{\'a}}1 {é}{{\'e}}1 {í}{{\'i}}1 {ó}{{\'o}}1 {ú}{{\'u}}1 {ñ}{{\~n}}1 {ü}{{\"u}}1
}
\lstdefinestyle{yamlstyle}{
    language=bash, % no hay lenguaje YAML nativo, se usa bash como base
    backgroundcolor=\color{backcolour},
    commentstyle=\color{codegreen},
    keywordstyle=\color{blue},
    numberstyle=\tiny\color{codegray},
    stringstyle=\color{codepurple},
    basicstyle=\ttfamily\scriptsize,
    breaklines=true,
    frame=single,
    showstringspaces=false,
    literate={á}{{\'a}}1 {é}{{\'e}}1 {í}{{\'i}}1 {ó}{{\'o}}1 {ú}{{\'u}}1 {ñ}{{\~n}}1 {ü}{{\"u}}1
}
\lstset{style=bashstyle} % Por defecto bash

% --- Pie de pagina personalizado ---
\makeatletter
\setbeamertemplate{footline}{
  \leavevmode%
  \hbox{%
  \begin{beamercolorbox}[wd=.333333\paperwidth,ht=2.25ex,dp=1ex,center]{author in head/foot}%
    \usebeamerfont{author in head/foot}\insertshortauthor
  \end{beamercolorbox}%
  \begin{beamercolorbox}[wd=.333333\paperwidth,ht=2.25ex,dp=1ex,center]{title in head/foot}%
    \usebeamerfont{title in head/foot}Sesión 13
  \end{beamercolorbox}%
  \begin{beamercolorbox}[wd=.333333\paperwidth,ht=2.25ex,dp=1ex,right]{date in head/foot}%
    \usebeamerfont{date in head/foot}BBDD \hfill \insertframenumber/\inserttotalframenumber \hspace*{0.3cm}
  \end{beamercolorbox}}%
  \vskip0pt%
}
\makeatother

% --- Datos de la portada ---
\title[SESIÓN 13 BBDD]{\textbf{Sesión 13}}
\subtitle{Despliegue y Alta Disponibilidad}
\author[Prof. C. Sánchez]{Carlos César Sánchez Coronel}
\institute{\textbf{BASES DE DATOS}}
\date{2026}

\setbeamertemplate{navigation symbols}{}

\AtBeginSection[]{
  \begin{frame}
    \frametitle{Contenido de la sección}
    \tableofcontents[currentsection]
  \end{frame}
}

\begin{document}

\begin{frame}
    \titlepage
\end{frame}

\begin{frame}{Contenido general}
    \tableofcontents
\end{frame}

% ============================================================================
% SECCIÓN 1: INTRODUCCIÓN
% ============================================================================
\section{Introducción}

\begin{frame}{Objetivos de la sesión}
    \begin{itemize}
        \item Comprender la importancia de la contenedorización para bases de datos.
        \item Aprender a desplegar SGBD en Docker y configurar persistencia.
        \item Conocer los conceptos básicos de Kubernetes aplicados a bases de datos.
        \item Diferenciar estrategias de alta disponibilidad: réplicas de lectura, failover, clustering.
        \item Configurar replicación en motores populares (PostgreSQL, MySQL).
        \item Entender el balanceo de carga para bases de datos.
        \item Analizar opciones en la nube (RDS, Cloud SQL) y comparar con on-premise.
        \item Planificar y ejecutar migraciones de bases de datos.
        \item Realizar ejercicios prácticos paso a paso.
    \end{itemize}
\end{frame}

% ============================================================================
% SECCIÓN 2: CONTENEDORES Y BASES DE DATOS
% ============================================================================
\section{Contenedores y Bases de Datos}

\begin{frame}{¿Por qué contenedores?}
    \begin{itemize}
        \item Entornos consistentes (desarrollo, prueba, producción).
        \item Rápido aprovisionamiento y destrucción.
        \item Aislamiento de recursos.
        \item Integración con CI/CD.
    \end{itemize}
\end{frame}

\begin{frame}[fragile]{Ejemplo: PostgreSQL en Docker}
    \begin{lstlisting}[style=bashstyle]
# Descargar imagen
docker pull postgres:15

# Ejecutar contenedor con persistencia
docker run -d \
  --name mi-postgres \
  -e POSTGRES_PASSWORD=mi_password \
  -e POSTGRES_USER=usuario \
  -e POSTGRES_DB=mi_db \
  -v pgdata:/var/lib/postgresql/data \
  -p 5432:5432 \
  postgres:15
    \end{lstlisting}
    \begin{itemize}
        \item \texttt{-v pgdata:/var/lib/postgresql/data}: volumen nombrado para persistencia.
        \item Variables de entorno configuran usuario, contraseña y base de datos inicial.
    \end{itemize}
\end{frame}

\begin{frame}{Volúmenes en Docker}
    \begin{description}
        \item[Volumen nombrado] Gestionado por Docker, ubicado en \texttt{/var/lib/docker/volumes/}. Recomendado para producción.
        \item[Bind mount] Monta un directorio del host. Útil en desarrollo.
    \end{description}
    \begin{exampleblock}{Comandos}
        \texttt{docker volume create pgdata}\\
        \texttt{docker run -v pgdata:/var/lib/postgresql/data ...}
    \end{exampleblock}
\end{frame}

\begin{frame}[fragile]{Redes en Docker}
    Para comunicación entre contenedores:
    \begin{lstlisting}[style=bashstyle]
docker network create mi-red
docker run -d --name postgres --network mi-red -e POSTGRES_PASSWORD=... postgres
docker run -d --name app --network mi-red -p 8080:8080 mi-app
    \end{lstlisting}
    La aplicación accede a la BD usando el hostname \texttt{postgres}.
\end{frame}

\begin{frame}{Consideraciones de Rendimiento en Producción}
    \begin{itemize}
        \item Usar volúmenes rápidos (SSD).
        \item Limitar recursos con \texttt{--memory} y \texttt{--cpus}.
        \item Ajustar parámetros del kernel (ej. \texttt{shared\_buffers}).
        \item Configurar rotación de logs.
    \end{itemize}
\end{frame}

% ============================================================================
% SECCIÓN 3: ORQUESTACIÓN CON KUBERNETES
% ============================================================================
\section{Orquestación con Kubernetes}

\begin{frame}{Kubernetes para Bases de Datos}
    \begin{itemize}
        \item Despliegues declarativos (Deployments, StatefulSets).
        \item Auto-escalado.
        \item Balanceo de carga interno (Services).
        \item Auto-reparación (reinicio de contenedores fallidos).
    \end{itemize}
\end{frame}

\begin{frame}{StatefulSet vs Deployment}
    \begin{description}
        \item[Deployment] Para aplicaciones sin estado. Pods intercambiables.
        \item[StatefulSet] Para aplicaciones con estado. Identidades estables (nombres como \texttt{postgres-0}), orden en despliegue/escalado, y almacenamiento persistente dedicado.
    \end{description}
    \alert{StatefulSet es el adecuado para bases de datos.}
\end{frame}

\begin{frame}[fragile]{Ejemplo: StatefulSet para PostgreSQL}
    \begin{lstlisting}[style=yamlstyle]
apiVersion: apps/v1
kind: StatefulSet
metadata:
  name: postgres
spec:
  serviceName: postgres
  replicas: 1
  selector:
    matchLabels:
      app: postgres
  template:
    metadata:
      labels:
        app: postgres
    spec:
      containers:
      - name: postgres
        image: postgres:15
        env:
        - name: POSTGRES_PASSWORD
          value: "mi_password"
        - name: POSTGRES_USER
          value: "usuario"
        - name: POSTGRES_DB
          value: "mi_db"
        ports:
        - containerPort: 5432
        volumeMounts:
        - name: pgdata
          mountPath: /var/lib/postgresql/data
  volumeClaimTemplates:
  - metadata:
      name: pgdata
    spec:
      accessModes: [ "ReadWriteOnce" ]
      resources:
        requests:
          storage: 10Gi
    \end{lstlisting}
\end{frame}

\begin{frame}[fragile]{Service para acceder a PostgreSQL}
    \begin{lstlisting}[style=yamlstyle]
apiVersion: v1
kind: Service
metadata:
  name: postgres
spec:
  selector:
    app: postgres
  ports:
    - port: 5432
      targetPort: 5432
  clusterIP: None  # headless para descubrimiento
    \end{lstlisting}
\end{frame}

\begin{frame}{Operadores para Bases de Datos en Kubernetes}
    Extienden Kubernetes para tareas complejas: backups, failover, upgrades.
    \begin{itemize}
        \item Crunchy PostgreSQL Operator
        \item Zalando PostgreSQL Operator
        \item Oracle MySQL Operator
    \end{itemize}
\end{frame}

% ============================================================================
% SECCIÓN 4: ESTRATEGIAS DE ALTA DISPONIBILIDAD (HA)
% ============================================================================
\section{Estrategias de Alta Disponibilidad}

\begin{frame}{Conceptos: RTO y RPO}
    \begin{description}
        \item[RTO] Recovery Time Objective: tiempo máximo para restaurar servicio.
        \item[RPO] Recovery Point Objective: pérdida máxima de datos aceptable.
    \end{description}
    Las estrategias HA buscan minimizar RTO y RPO.
\end{frame}

\begin{frame}{Réplicas de Lectura}
    \begin{itemize}
        \item Copias de la base principal que reciben cambios y aceptan lecturas.
        \item Beneficios: escalabilidad de lecturas, aislamiento de cargas, mayor disponibilidad.
        \item Implementaciones: PostgreSQL streaming replication, MySQL replication, MongoDB replica sets.
    \end{itemize}
\end{frame}

\begin{frame}{Failover}
    Conmutación a nodo secundario ante fallo del primario.
    \begin{description}
        \item[Manual] Administrador promueve la réplica.
        \item[Automático] Sistema de monitoreo ejecuta promoción (Patroni, repmgr, Orchestrator).
    \end{description}
    \alert{Riesgo de split-brain: usar quorum (etcd, ZooKeeper) y fencing.}
\end{frame}

\begin{frame}{Clustering}
    \begin{description}
        \item[Activo-Pasivo] Un nodo activo, otros en espera. Ej: PostgreSQL con Patroni.
        \item[Activo-Activo] Todos los nodos aceptan escrituras, con resolución de conflictos. Ej: MySQL Cluster, Oracle RAC, Cassandra.
    \end{description}
\end{frame}

\begin{frame}{Balanceo de Carga}
    Distribuye tráfico entre varios servidores.
    \begin{itemize}
        \item \textbf{HAProxy}: balanceo TCP genérico.
        \item \textbf{ProxySQL}: balanceo y enrutamiento para MySQL.
        \item \textbf{pgpool-II}: pooling y balanceo para PostgreSQL.
    \end{itemize}
\end{frame}

% ============================================================================
% SECCIÓN 5: REPLICACIÓN EN MOTORES ESPECÍFICOS
% ============================================================================
\section{Replicación en Motores Específicos}

\begin{frame}[fragile]{PostgreSQL: Streaming Replication}
    \begin{block}{Primario}
        \begin{lstlisting}[style=sqlstyle]
CREATE USER replicator WITH REPLICATION ENCRYPTED PASSWORD 'pass';
        \end{lstlisting}
        En \texttt{postgresql.conf}:
        \begin{lstlisting}[style=bashstyle]
wal_level = replica
max_wal_senders = 5
wal_keep_size = 1GB
        \end{lstlisting}
        En \texttt{pg\_hba.conf}:
        \begin{lstlisting}[style=bashstyle]
host replication replicator 192.168.1.0/24 md5
        \end{lstlisting}
    \end{block}
    \begin{block}{Secundario}
        \begin{lstlisting}[style=bashstyle]
pg_basebackup -h primario -D /var/lib/postgresql/data -U replicator -v -P --wal-method=stream
touch /var/lib/postgresql/data/standby.signal
echo "primary_conninfo = 'host=primario user=replicator password=pass'" >> postgresql.auto.conf
        \end{lstlisting}
    \end{block}
\end{frame}

\begin{frame}{MySQL: Group Replication}
    \begin{itemize}
        \item Replicación multi-maestro con detección automática de fallos.
        \item Configuración mediante MySQL Shell (InnoDB Cluster).
        \item Proporciona consistencia fuerte y failover automático.
    \end{itemize}
\end{frame}

\begin{frame}{MongoDB: Replica Sets}
    \begin{itemize}
        \item Grupo de servidores mongod con un primario y varios secundarios.
        \item Failover automático: si el primario falla, los secundarios eligen uno nuevo.
        \item Configuración: iniciar cada nodo con \texttt{--replSet}, luego \texttt{rs.initiate()} y \texttt{rs.add()}.
    \end{itemize}
\end{frame}

\begin{frame}{Oracle: Data Guard}
    \begin{itemize}
        \item Mantiene bases de datos standby (físicas o lógicas).
        \item Puede ser síncrona (protección máxima) o asíncrona.
        \item Fast-Start Failover para conmutación automática.
    \end{itemize}
\end{frame}

\begin{frame}{SQL Server: Always On Availability Groups}
    \begin{itemize}
        \item Grupos de disponibilidad con réplicas síncronas o asíncronas.
        \item Hasta 9 réplicas, failover automático con cluster de Windows.
        \item Réplicas legibles.
    \end{itemize}
\end{frame}

% ============================================================================
% SECCIÓN 6: BALANCEO DE CARGA PRÁCTICO
% ============================================================================
\section{Balanceo de Carga Práctico}

\begin{frame}[fragile]{HAProxy para PostgreSQL}
    Configuración en \texttt{/etc/haproxy/haproxy.cfg}:
    \begin{lstlisting}[style=bashstyle]
global
    log stdout format raw local0
defaults
    log global
    option tcplog
    timeout connect 5s
    timeout client 1m
    timeout server 1m
frontend pgsql
    bind *:5432
    default_backend pgsql_backend
backend pgsql_backend
    balance roundrobin
    option pgsql-check user haproxy
    server pg1 192.168.1.10:5432 check
    server pg2 192.168.1.11:5432 check
    \end{lstlisting}
    Crear usuario \texttt{haproxy} en PostgreSQL para health checks.
\end{frame}

\begin{frame}{ProxySQL para MySQL}
    \begin{itemize}
        \item Enrutamiento de consultas (lecturas a réplicas, escrituras al primario).
        \item Pooling de conexiones, caché de consultas, control de tráfico.
        \item Configuración mediante tablas en memoria.
    \end{itemize}
\end{frame}

\begin{frame}{pgpool-II para PostgreSQL}
    \begin{itemize}
        \item Middleware con pooling de conexiones, balanceo de carga, replicación y failover.
        \item Se sitúa entre la aplicación y los servidores PostgreSQL.
    \end{itemize}
\end{frame}

% ============================================================================
% SECCIÓN 7: DESPLIEGUE EN NUBE VS ON-PREMISE Y MIGRACIÓN
% ============================================================================
\section{Despliegue en Nube vs On-Premise y Migración}

\begin{frame}{Servicios Gestionados (PaaS)}
    \begin{description}
        \item[AWS RDS] Soporta múltiples motores, replicación, backups, failover multi-AZ.
        \item[Google Cloud SQL] PostgreSQL, MySQL, SQL Server con alta disponibilidad.
        \item[Azure SQL Database] Base de datos como servicio escalable.
    \end{description}
    \textbf{Ventajas:} menor overhead, escalado sencillo, HA integrada. \\
    \textbf{Desventajas:} menos control, vendor lock-in, costos variables.
\end{frame}

\begin{frame}{On-Premise}
    \textbf{Ventajas:} control total, cumplimiento normativo, costos fijos. \\
    \textbf{Desventajas:} inversión inicial alta, mantenimiento complejo, escalado limitado.
\end{frame}

\begin{frame}{Arquitectura Híbrida}
    Combinación de on-premise y nube. Ej: datos sensibles on-premise, data lake en la nube para análisis, replicación para disaster recovery.
\end{frame}

\begin{frame}[fragile]{Migración a la Nube: Ejemplo con pg\_dump}
    \begin{lstlisting}[style=bashstyle]
# Dump local
pg_dump -U usuario -h localhost mi_db > mi_db.sql

# Restaurar en RDS (previamente crear instancia)
psql -U usuario -h mi-instance.rds.amazonaws.com -d mi_db < mi_db.sql
    \end{lstlisting}
    Para migración sin downtime: usar AWS DMS con replicación continua (CDC).
\end{frame}

\begin{frame}{AWS Database Migration Service (DMS)}
    \begin{itemize}
        \item Servicio gestionado para migraciones con mínimo downtime.
        \item Soporta replicación continua desde el origen.
        \item Ideal para grandes volúmenes o migraciones heterogéneas (ej. Oracle a PostgreSQL).
    \end{itemize}
\end{frame}

% ============================================================================
% SECCIÓN 8: EJERCICIOS RESUELTOS
% ============================================================================
\section{Ejercicios Resueltos}

\begin{frame}[fragile]{Ejercicio 1: PostgreSQL con Docker y persistencia}
    \textbf{Enunciado:} Crear contenedor PostgreSQL con volumen persistente, crear tabla, insertar datos, detener y eliminar contenedor, luego crear nuevo con mismo volumen y verificar persistencia.
    \begin{block}{Solución}
        \begin{lstlisting}[style=bashstyle]
docker volume create pgdata
docker run -d --name postgres1 -e POSTGRES_PASSWORD=admin -e POSTGRES_USER=admin -e POSTGRES_DB=testdb -v pgdata:/var/lib/postgresql/data -p 5432:5432 postgres:15
docker exec -it postgres1 psql -U admin -d testdb
# En psql:
CREATE TABLE personas (id SERIAL, nombre TEXT);
INSERT INTO personas (nombre) VALUES ('Ana'), ('Luis');
SELECT * FROM personas;
\q
docker stop postgres1 && docker rm postgres1
docker run -d --name postgres2 -e POSTGRES_PASSWORD=admin -e POSTGRES_USER=admin -e POSTGRES_DB=testdb -v pgdata:/var/lib/postgresql/data -p 5432:5432 postgres:15
docker exec -it postgres2 psql -U admin -d testdb -c "SELECT * FROM personas;"
        \end{lstlisting}
    \end{block}
\end{frame}

\begin{frame}[fragile]{Ejercicio 2: Configurar replicación PostgreSQL con Docker}
    \textbf{Enunciado:} Levantar dos contenedores PostgreSQL (primario y réplica) en una misma máquina y configurar streaming replication asíncrona.
    \begin{block}{Solución (resumida)}
        \begin{lstlisting}[style=bashstyle]
docker network create pg-net
# Primario
docker run -d --name pg-primary --network pg-net -e POSTGRES_PASSWORD=admin -e POSTGRES_USER=admin -e POSTGRES_DB=testdb -v pg-primary-data:/var/lib/postgresql/data postgres:15
# Configurar primario (wal_level, max_wal_senders, pg_hba) y crear usuario replicator
# Réplica: hacer basebackup y configurar standby
        \end{lstlisting}
        Ver detalles en el solucionario.
    \end{block}
\end{frame}

\begin{frame}[fragile]{Ejercicio 3: Failover manual en PostgreSQL}
    \textbf{Enunciado:} Simular caída del primario y promover la réplica.
    \begin{block}{Solución}
        \begin{lstlisting}[style=bashstyle]
docker stop pg-primary
docker exec -it pg-replica pg_ctl promote -D /var/lib/postgresql/data
# Ahora pg-replica es primario
        \end{lstlisting}
    \end{block}
\end{frame}

\begin{frame}[fragile]{Ejercicio 4: Migrar MySQL local a AWS RDS (simulado)}
    \textbf{Enunciado:} Migrar base de datos MySQL desde contenedor local a otro contenedor que simula RDS.
    \begin{block}{Solución}
        \begin{lstlisting}[style=bashstyle]
# Origen
docker run -d --name mysql-origen -e MYSQL_ROOT_PASSWORD=root -e MYSQL_DATABASE=ventas -p 3306:3306 mysql:8
# Destino
docker run -d --name mysql-destino -e MYSQL_ROOT_PASSWORD=root -e MYSQL_DATABASE=ventas -p 3307:3306 mysql:8
# Dump y restore
mysqldump -h127.0.0.1 -P3306 -uroot -proot ventas > ventas.sql
mysql -h127.0.0.1 -P3307 -uroot -proot ventas < ventas.sql
        \end{lstlisting}
    \end{block}
\end{frame}

\begin{frame}[fragile]{Ejercicio 5: Configurar HAProxy para balancear entre dos réplicas PostgreSQL}
    \textbf{Enunciado:} Usar HAProxy para balancear conexiones de solo lectura entre dos réplicas.
    \begin{block}{Solución}
        \begin{lstlisting}[style=bashstyle]
# Crear red y conectar réplicas
docker network create ha-net
# Ejecutar HAProxy con configuración (ver diapositiva anterior)
# En cada réplica, crear usuario haproxy
docker exec -it pg-replica1 psql -U admin -c "CREATE USER haproxy;"
# Probar conexión
psql -h127.0.0.1 -p5432 -U admin -d testdb -c "SELECT inet_server_addr();"
        \end{lstlisting}
    \end{block}
\end{frame}

% ============================================================================
% SECCIÓN 9: GLOSARIO
% ============================================================================
\section{Glosario}

\begin{frame}{Términos Clave}
    \begin{description}
        \item[Contenedor] Entorno aislado para ejecutar aplicaciones.
        \item[Docker] Plataforma de contenedores.
        \item[Volumen] Mecanismo de persistencia en Docker.
        \item[Kubernetes] Orquestador de contenedores.
        \item[StatefulSet] Recurso de Kubernetes para aplicaciones con estado.
        \item[Replicación] Copia de datos de un servidor a otro.
        \item[Réplica de lectura] Copia que acepta consultas de solo lectura.
        \item[Failover] Conmutación a nodo secundario ante fallo.
        \item[Split-brain] Dos nodos creyéndose primarios.
        \item[RTO] Recovery Time Objective.
        \item[RPO] Recovery Point Objective.
        \item[HAProxy] Balanceador TCP/HTTP.
        \item[ProxySQL] Proxy para MySQL.
        \item[pgpool-II] Middleware para PostgreSQL.
        \item[AWS RDS] Servicio de base de datos relacional gestionado.
    \end{description}
\end{frame}

% --- Diapositiva final ---
\begin{frame}
    \centering
    \Huge \textbf{¿Preguntas?}
    
    \vspace{1cm}
    \large ¡Gracias por su atención!
\end{frame}

\end{document}